\documentclass{article}

\usepackage{booktabs}
\usepackage{tabularx}

\title{Development Plan\\\progname}

\author{\authname}

\date{}

\input{../Comments.text}
\input{../Common.text}

\begin{document}

\maketitle

\begin{table}[hp]
\caption{Revision History} \label{TblRevisionHistory}
\begin{tabularx}{\textwidth}{llX}
\toprule
\textbf{Date} & \textbf{Developer(s)} & \textbf{Change}\\
\midrule
Date1 & Name(s) & Description of changes\\
Date2 & Name(s) & Description of changes\\
... & ... & ...\\
\bottomrule
\end{tabularx}
\end{table}

\newpage{}

\wss{Put your introductory blurb here.  Often the blurb is a brief roadmap of
what is contained in the report.}

\wss{Additional information on the development plan can be found in the
\href{https://gitlab.cas.mcmaster.ca/courses/capstone/-/blob/main/Lectures/L02b_POCAndDevPlan/POCAndDevPlan.pdf?ref_type=heads}
{lecture slides}.}

\section{Confidential Information}

There is no confidential information to protect.

\section{IP to Protect}

There is no confidential information to protect

\section{Copyright License}

Lake is licensed under the GNU General Public License Version 3.
A copy of the GNU GPL license is available in the Lake repository
\href{https://github.com/capstoneCEGJM/Lake/blob/main/LICENSE}
{here}.

\section{Team Meeting Plan}

CinqGarçons will meet at minimum weekly every Friday at 2:40PM EST via Discord. Additional ad-hoc meetings
will be scheduled as the need arises. Some ad-hoc meetings may be in-person, likely after/before
lecture.

\vspace{5mm}

Each week, team members will take turns being the chair and leading the meeting. The chair will be
responsible for bringing a formal agenda to each meeting. The chair will also be responsible for
introducing the agenda to the team, and ensuring the agenda points are addressed during the meeting.
Ad-hoc meetings may only have an informal unwritten agenda due to their potentially last-minute nature.

\vspace{5mm}

The order in which members will fill the role of meeting chair are as follows:
\begin{itemize}
	\item Colin Chambachan
	\item Jinwoo Hong
	\item Matthew Nesbitt
	\item Gary Qin
	\item Ethan Walsh
\end{itemize}

\section{Team Communication Plan}

Formal text-based team communications will be facilitated through a Discord server, with additional
communication happening verbally during meetings.

\vspace{5mm}

Project design and repository related communications may also occur on a pull request during the
review phase, or in the comments of a related GitHub issue. All formal action items shall have an
associated GitHub issue to track the item's progress. Pull requests should be associated with an
issue, and may resolve said Issue. Issues will be the primary method of work assignment and work
splitting within the team, as each issue shall have one or more assigned members. Issues may also be
created as a discussion thread for particular design decisions that require deliberation.

\section{Team Member Roles}

Team members may fill roles dynamically throughout the year as best fits the situation. Some roles,
however, are assigned permanent members. Members may fill more than one role, and roles may be
filled by more than one member. Identified member roles are as follows:
\begin{itemize}
	\item Liason: Ethan Walsh
	\item Meeting Chair: Dynamic
	\item Repo/Org Administrator: Matthew Nesbitt
	\item Technical Lead (Architecture): TBD
	\item Technical Lead (IPC): TBD
	\item Technical Lead (UI/UX): TBD
	\item Technical Lead (ML): TBD
	\item Technical Lead (Performance): TBD
\end{itemize}

\section{Workflow Plan}

All contributions to production are required to go through a pull request, and aquire approval from
at minimum two reviewers. When working on a contribution, using either the branching workflow or
forking workflow is acceptable, and it is up to the individual contributor to decide which workflow
to use.

\subsection{Contribution Flow}
\begin{enumerate}
	\item For branching workflow:
		\begin{enumerate}
			\item Create a new working branch in the main repository. The branch name should be
				descriptive, and should be prefixed with the contributor's name (i.e.
				\verb|<name>/<branch>|)
			\item Commit desired changes to the working branch. Commit names should be descriptive
				but turse, and should only make one individual change.
		\end{enumerate}
	\item For forking workflow:
		\begin{enumerate}
			\item Create a fork of the main repository. Only one fork per contributor is required.
			\item Create a new working branch in this fork. There are no restrictions on branch
				names.
			\item Commit desired changes to the working branch in the fork. Commit names should be
				descriptive but turse, and should only make one individual change.
		\end{enumerate}
	\item Open a pull request from the working branch to the \verb|main| branch in the main repository.
	\item Address any changes requested during review.
	\item After the pull request has recieved two approvals and passes CI, it can be merged.
	\item For branching workflow:
		\begin{enumerate}
			\item After a PR has been merged, the associated work branch should be deleted. This can
				be done directly from the PR page.
		\end{enumerate}
\end{enumerate}

\subsection{Issues}

Issues will be used to track action items, and occasionally for design decisions that require
deliberation. When creating an issue, a description should be included providing all the information
required to address the issue. In addition, the issue may be labeled with one or more relevant
tag(s) if applicable. The available tags are as follows: \verb|bug|, \verb|feature|, \verb|decision|,
\verb|improvement|, \verb|priority: high|, \verb|priority: low|. Pull requests should be associated
with an issue, and may close said issue by including a close tag.

\vspace{5mm}

Issues will also be created (in seperate projects) for each lecture, as well as each TA meeting.
Issue templates will be created for both lecture and TA meeting issues to make the creation
of these issues easier and faster.

\subsection{CI/CD}

CI/CD will utilize GitHub Actions, will be required to pass for a PR to be merged, and will
consist of the following stages:
\begin{itemize}
	\item PDF building and publishing
	\item Code linting
	\item Unit testing
	\item Integration testing
	\item System testing
	\item Performance testing
\end{itemize}

\section{Project Decomposition and Scheduling}

\begin{itemize}
  \item How will you be using GitHub projects?
  \item Include a link to your GitHub project
\end{itemize}

\wss{How will the project be scheduled?  This is the big picture schedule, not
details. You will need to reproduce information that is in the course outline
for deadlines.}

\section{Proof of Concept Demonstration Plan}

What is the main risk, or risks, for the success of your project?  What will you
demonstrate during your proof of concept demonstration to convince yourself that
you will be able to overcome this risk?

\section{Expected Technology}

\wss{What programming language or languages do you expect to use?  What external
libraries?  What frameworks?  What technologies.  Are there major components of
the implementation that you expect you will implement, despite the existence of
libraries that provide the required functionality.  For projects with machine
learning, will you use pre-trained models, or be training your own model?  }

\wss{The implementation decisions can, and likely will, change over the course
of the project.  The initial documentation should be written in an abstract way;
it should be agnostic of the implementation choices, unless the implementation
choices are project constraints.  However, recording our initial thoughts on
implementation helps understand the challenge level and feasibility of a
project.  It may also help with early identification of areas where project
members will need to augment their training.}

Topics to discuss include the following:

\begin{itemize}
\item Specific programming language
\item Specific libraries
\item Pre-trained models
\item Specific linter tool (if appropriate)
\item Specific unit testing framework
\item Investigation of code coverage measuring tools
\item Specific plans for Continuous Integration (CI), or an explanation that CI
  is not being done
\item Specific performance measuring tools (like Valgrind), if
  appropriate
\item Tools you will likely be using?
\end{itemize}

\wss{git, GitHub and GitHub projects should be part of your technology.}

\section{Use of AI and LLMs}

\wss{This section is a continuation of the previous section.  The role of LLMs
in your project as important enough to warrant their own section.  
How will your team be using LLMs for code and documentation development?  
What tools will you be using?  How will you verify the output of your LLMs?}

\section{Coding Standard}

\wss{What coding standard will you adopt?}

\newpage{}

\section{Appendix --- Generative AI Use Disclosure}

\wss{ChatGPT (version 5) was used to brainstorm ideas for the template for this disclose section.}

\wss{This section is about the use of AI to write this document, not your
planned use of AI for your project.  You discuss that topic in one of the
sections in the body of this report.}

\subsection{AI Tools Used}

\wss{Give the name of the AI tool(s) used and the version.}

\subsection{How and Where Used}

\wss{Explain what the tool(s) were used for.  Examples include: brainstorming,
explaining a technical concept, reviewing the document, generating or modifying
text, generating or modifying diagrams, identifying errors or omissions or
inconsistencies.}

\wss{For each of the uses, identify which parts of the document were affected.}

\wss{You do not need to report every interaction, but you should disclose the 
uses that materially contributed to your work}

\subsection{Verification of AI-Generated Content}

\wss{\begin{itemize}
    \item Describe how the team checked AI-generated or AI-assisted material
    before including it in the document.
    \item In particular, verify factual claims, technical assertions,
    requirements, calculations, references, and other information for which an
    AI system may produce incorrect or fabricated results.
    \item \textbf{The team is responsible for the accuracy, completeness,
consistency, and appropriateness of the submitted document, regardless of
whether material was generated by a human or an AI system.}
\end{itemize}}

\newpage{}

\section*{Appendix --- Reflection}

\wss{Not required for CAS 741}

\input{../Reflection.text}

\begin{enumerate}
    \item Why is it important to create a development plan prior to starting the
    project?
    \item In your opinion, what are the advantages and disadvantages of using
    CI/CD?
    \item What disagreements did your group have in this deliverable, if any,
    and how did you resolve them?
\end{enumerate}

\newpage{}

\section*{Appendix --- Team Charter}

\wss{borrows from
\href{https://engineering.up.edu/industry_partnerships/files/team-charter.pdf}
{University of Portland Team Charter}}

\subsection*{External Goals}

\wss{What are your team's external goals for this project? These are not the
goals related to the functionality or quality fo the project.  These are the
goals on what the team wishes to achieve with the project.  Potential goals are
to win a prize at the Capstone EXPO, or to have something to talk about in
interviews, or to get an A+, etc.}

\subsection*{Attendance}

\subsubsection*{Expectations}

\wss{What are your team's expectations regarding meeting attendance (being on
time, leaving early, missing meetings, etc.)?}

\subsubsection*{Acceptable Excuse}

\wss{What constitutes an acceptable excuse for missing a meeting or a deadline?
What types of excuses will not be considered acceptable?}

\subsubsection*{In Case of Emergency}

\wss{What process will team members follow if they have an emergency and cannot
attend a team meeting or complete their individual work promised for a team
deliverable?}

\subsection*{Accountability and Teamwork}

\subsubsection*{Quality} 

\wss{What are your team's expectations regarding the quality
of team members' preparation for team meetings and the quality of the
deliverables that members bring to the team?}

\subsubsection*{Attitude}

\wss{What are your team's expectations regarding team members' ideas,
interactions with the team, cooperation, attitudes, and anything else regarding
team member contributions?  Do you want to introduce a code of conduct?  Do you
want a conflict resolution plan?  Can adopt existing codes of conduct.}

\subsubsection*{Stay on Track}

\wss{What methods will be used to keep the team on track? How will your team
ensure that members contribute as expected to the team and that the team
performs as expected? How will your team reward members who do well and manage
members whose performance is below expectations?  What are the consequences for
someone not contributing their fair share?}

\wss{You may wish to use the project management metrics collected for the TA and
instructor for this.}

\wss{You can set target metrics for attendance, commits, etc.  What are the
consequences if someone doesn't hit their targets?  Do they need to bring the
coffee to the next team meeting?  Does the team need to make an appointment with
their TA, or the instructor?  Are there incentives for reaching targets early?}

\subsubsection*{Team Building}

\wss{How will you build team cohesion (fun time, group rituals, etc.)? }

\subsubsection*{Decision Making} 

\wss{How will you make decisions in your group? Consensus?  Vote? How will you
handle disagreements? }

\end{document}
